\section{问题重述}\label{sec:problem}

\subsection{竞赛背景与预测任务}
天池精准医疗大赛初赛以临床数据和体检指标为输入，以个体血糖值为预测目标。尽管赛事总体背景涉及糖尿病及遗传风险，本文所用数据和模型均对应初赛的连续血糖回归任务。研究不包含遗传变量建模，也不将血糖回归结果转换为未经验证的糖尿病分类或诊断结论。

数据按记录行组织，每行表示一个个体，首列为ID。训练表包含42列，最后一列为目标血糖值；本地A、B特征表各包含41列，不含目标列。字段涉及性别、年龄、体检日期及多类检验指标，部分检验值缺失。研究需要在有限样本及不完整输入条件下，为每条待预测记录给出连续数值。

\subsection{输出与评价要求}
设评价集合为$E$，样本数为$n_E$，第$i$条记录的真实值与预测值分别为$y_i$和$\hat y_i$。竞赛以均方误差评价预测结果：
\begin{equation}
\operatorname{MSE}_E=\frac{1}{n_E}\sum_{i\in E}(y_i-\hat y_i)^2.
\label{eq:mse}
\end{equation}
误差越低，说明该评价集合上的平方损失越小。初赛提交要求预测值保留小数点后三位。本文报告的是既有实验基于未舍入预测计算的指标，三位小数属于最终提交的数值格式要求，不能将两种计算口径混称为同一提交成绩。

题目还提出高精度、高效和解释性要求。本文以同协议预测误差考察精度，以实际搜索预算说明计算范围，并通过线性探索、特征消融、融合组成与误差诊断分析模型行为。现有研究未包含统一硬件下的效率比较或最终融合的个体预测解释验证，因此相关讨论限定于已有证据。
